% Generated draft: requires mathematical and source review.
\documentclass[11pt]{article}
\usepackage[margin=25mm]{geometry}
\usepackage{fontspec}
\setmainfont{Latin Modern Roman}
\usepackage{amsmath,amssymb,mathtools,unicode-math}
\setmathfont{Latin Modern Math}
\usepackage{enumitem,needspace,xurl,hyperref}
\setlength{\parindent}{0pt}
\setlength{\parskip}{5pt}
\setlength{\emergencystretch}{3em}
\newcommand{\N}{\mathbb{N}}
\newcommand{\Z}{\mathbb{Z}}
\newcommand{\Q}{\mathbb{Q}}
\newcommand{\R}{\mathbb{R}}
\newcommand{\C}{\mathbb{C}}
\newcommand{\F}{\mathbb{F}}
\newcommand{\A}{\mathbb{A}}
\newcommand{\PP}{\mathbb{P}}
\newcommand{\E}{\mathbb{E}}
\newcommand{\eps}{\varepsilon}
\newcommand{\dd}{\,\mathrm{d}}
\begin{document}
\section*{UnsolvedMath ID: 3700016}
\textbf{Source problem code:} catalogue entry
\subsection{Definitions and mathematical statement}
Set $\omega=\exp(2\pi i/5)$ and let
\[
\mathcal S=\{f:\C\to\C\text{ entire}:\ 
 f(\omega z)f(\omega^{-1}z)=f(z)-1\text{ for every }z\in\C\}.
\]
Under the usual meaning of rotation of the independent variable, the literal uniqueness question is whether
\[
\forall f,g\in\mathcal S\ \exists\theta\in\R\ \forall z\in\C,
\qquad g(z)=f(e^{i\theta}z).
\]
This is uniqueness of a rotation orbit, not an assertion of a specified growth order or normalization. If ``unique'' also includes existence, one must additionally require $\mathcal S\ne\varnothing$.

The supplied question states no nonconstancy condition, normalization, growth condition, or restriction on the rotation angle. Constants obey $c^2=c-1$, so the literal class includes two different constant solutions, which rotations cannot identify. Also, the equation is invariant under $f(z)\mapsto f(az)$ for every complex $a$, and nonzero dilation is a broader equivalence than rotation. These features require clarification of the intended normalized solution class or equivalence relation before the source's open-uniqueness interpretation can be adopted.

The research summary also cites the distinct equation
\[
f(\lambda)+f(\omega\lambda)f(\omega^{-1}\lambda)=1.
\]
Its sign convention is not the displayed source equation; no identification between their solution classes is assumed here.
\subsection{Short English statement}
For the stated fifth-root functional equation, does every entire solution differ from every other only by rotating the argument? The literal question includes inequivalent constant solutions and gives no normalization to control dilation, so its intended open formulation needs clarification.
\subsection{Sources}
\begin{itemize}
\item Local UnsolvedMath registry, permanent ID 3700016, exact source code catalogue entry, ``A fifth-root functional equation''; proposed by Alexandre Eremenko and cited proposers, proposed year 2017. Registry status open, classification OPEN-TRIAGE.
\item Alexandre Eremenko, \emph{My favorite unsolved problems}, ``Radially distributed values'', second question, \url{https://www.math.purdue.edu/~eremenko/uns1.html}. Supplied provenance: accessed July 29, 2026, extraction pdf-text; retained after applying the author index's explicit solved/counterexample updates through that date. No PDF was accessed for this draft.
\item The supplied summary quotes Eremenko's \emph{radial.pdf}: ``Very little is known about entire solutions of this functional equation [Sibuya]. Is an entire solution unique, up to rotation of $z$?'' It relates the fifth-fold symmetry to the $\alpha=2\pi/5$ radially distributed example, source code AMR-036-0015, PT-symmetric equations of Voros, and integrable models. No precise Voros or integrable-model reference is provided.
\item Y. Sibuya, reference supplied as \emph{Analysis} \textbf{8} (1998), 271--295, for the $m=3$ equation $f(\lambda)+f(\omega\lambda)f(\omega^{-1}\lambda)=1$; the summary also mentions earlier Sibuya--student work. No title or full bibliography is given, and the supplied volume/year and sign conventions need checking.
\item The supplied summary says no later published resolution of uniqueness up to rotation was verified. Provenance flags status and rights NEEDS\_REVIEW: public access was verified, but redistribution permission was not established. These are supplied assessments, not independent source or literature verification.
\end{itemize}
\end{document}
